\documentclass[10pt,letterpaper,twocolumn]{article}
\usepackage{smallsat2026}
\usepackage{amsmath}

% Paper-specific footer values for the shared SmallSat template.
\providecommand{\smallsatauthor}{Kankanala}
\providecommand{\smallsatconference}{41st Annual Small Satellite Conference}
\providecommand{\smallsatpaperid}{}
\lfoot{\footnotesize \smallsatauthor}
\rfoot{\footnotesize \smallsatconference}

% Figures are stored in figures/. The style package provides graphicx.
\graphicspath{{figures/}}

\title{CubeSat Design for Astrophysical ALP Observations using Skipper-CCDs}

\author{Roshni Kankanala}
\date{41st Annual Small Satellite Conference, 2027}

% Add the assigned code when available, e.g. SSC27-P1-00.
\renewcommand{\smallsatpaperid}{}

\begin{document}

\fancypagestyle{firstpage}{%
  \fancyhf{}%
  \rhead{\footnotesize \smallsatpaperid}%
  \lfoot{\footnotesize \smallsatauthor}%
  \cfoot{\footnotesize \thepage}%
  \rfoot{\footnotesize \smallsatconference}%
}

\thispagestyle{firstpage}

\twocolumn[
\maketitle
\begin{abstract}
% TODO: Add a concise summary of the problem, approach, and key findings.
\end{abstract}
\keywords{TODO: Add three to five keywords}
]

\section{Introduction}
% Student source: Draft 1 -- AE 298 RES, by Roshni Kankanala.
Dark matter searches probe candidate particles through their interactions with ordinary matter and their possible decay products. Underground detectors suppress cosmic-ray backgrounds, but the atmosphere and crust can attenuate strongly interacting dark matter before it reaches a detector. The atmosphere also absorbs astrophysical X-rays. Observations from low Earth orbit (LEO) provide access to these particle and photon signals\cite{alpine2025darkness}.

Charge-coupled devices (CCDs) measure charge generated by energy deposition in silicon. Absorbed X-rays produce electron--hole pairs, and the collected pixel charge provides an estimate of the deposited energy. Skipper-CCDs permit repeated, nondestructive measurements of the same charge packet. Averaging these measurements reduces electronic readout noise and supports searches for low-energy ionization and X-ray events\cite{alpine2025darkness}.

The published DarkNESS design is a 6U NanoSat intended to demonstrate Skipper-CCD operation in LEO and conduct two dark matter searches\cite{alpine2025darkness}.

\begin{enumerate}
  \item \textbf{Search for strongly interacting sub-GeV dark matter.}
  Dark matter scattering in silicon would produce low-energy ionization. Observations toward Cygnus target the expected dark matter wind associated with the Solar System's motion through the Galactic halo. Earth's changing shielding along the orbit would modulate the event rate.

  \item \textbf{Search for X-rays from dark matter decay.}
  Observations toward the Galactic Center target spectral lines from candidate dark matter decays. One example is the decay of a keV-scale sterile neutrino into an active neutrino and an X-ray photon. The published observing concept uses orbital umbra to limit solar contamination\cite{alpine2025darkness}.
\end{enumerate}

The ALP search considered here uses DarkNESS as a reference instrument for an additional science objective. Axion-like particles (ALPs) are hypothetical pseudoscalar particles whose production and interactions depend on the chosen particle model. The cosmological source channel supplies a diffuse ALP intensity from cosmological distances. The solar source channel supplies a directional ALP flux produced in the Sun. In both channels, ALPs could convert into X-ray photons in Earth's magnetic field, and the Skipper-CCDs would record the resulting photons\cite{yamamoto2020alp,davoudiasl2008gecosax}. Solar ALPs need not constitute the Galactic dark matter population.

The study addresses the following research question.

\begin{quote}
``Which astrophysical ALP source channels can a LEO X-ray observatory identify through conversion in Earth's magnetic field, and what instrument response and observing strategy would make those signals measurable against particle and diffuse X-ray backgrounds?''
\end{quote}

\begin{enumerate}
  \item How feasible is a cosmological ALP source channel for a LEO CubeSat using Skipper-CCDs, assuming conversion in Earth's magnetic field?
  \item How feasible is a solar ALP source channel for a LEO CubeSat using Skipper-CCDs, assuming conversion in Earth's magnetic field?
  \item How should the CubeSat's pointing strategy and observing constraints be defined for each candidate source channel?
  \item How should the CubeSat monitor backgrounds, using proxies, control exposures, and a particle monitor, to distinguish an ALP conversion signal?
\end{enumerate}

\section{Science and Background}
\subsection{Inverse Primakoff Effect and ALPs}
An ALP with a photon coupling can convert into a photon in a transverse magnetic field, approximately preserving its energy. The reference calculation assumes weak mixing, negligible absorption, and coherent conversion along the accepted ray\cite{yamamoto2020alp,davoudiasl2008gecosax}. In this limit, the probability is
\begin{equation}
\begin{aligned}
P_{a\gamma} &\simeq 2.45\times10^{-21}g_{10}^{2}
 \left(\frac{\mathcal{B}}{\mathrm{T\,m}}\right)^{2},\\
\mathcal{B} &\equiv\left|\int_0^L\mathbf{B}_{\perp}(s)\,ds\right|,
\qquad g_{10}\equiv\frac{g_{a\gamma}}{10^{-10}\,\mathrm{GeV}^{-1}}.
\end{aligned}
\label{eq:conversion}
\end{equation}
Here $g_{a\gamma}$ is the unknown photon coupling, $\mathbf{B}_{\perp}$ is the signed transverse field, and $L$ is the transparent conversion path. The coefficient includes the conversion of tesla and metres to natural units. Coherence requires a small accumulated phase. For a uniform medium, the phase mismatch is
\begin{equation}
qL=\frac{(m_a^2-m_\gamma^2)L}{2E},\qquad |qL|\ll1,
\qquad \hbar=c=1.
\label{eq:coherence}
\end{equation}
The ALP mass $m_a$, effective photon mass $m_\gamma$, and energy $E$ determine the mass reach. A complete ray calculation must include phase cancellation and atmospheric attenuation where these approximations fail.

\subsection{The Cosmological Channel}
The reference source assumes a nonrelativistic 7~keV dark matter parent with a lifetime much longer than the age of the Universe. Each decay produces two relativistic ALPs with 3.5~keV initial energy. Expansion redshifts this population into a diffuse continuum. Geomagnetic conversion then produces photons that can be absorbed in silicon. The continuum intensity per steradian follows from the decay rate and cosmological expansion\cite{yamamoto2020alp}.
\begin{equation}
\begin{aligned}
\frac{dI_{a,\mathrm{cos}}}{dE}
 &=\frac{c f\rho_{\mathrm{DM},0}}{2\pi m_\chi\tau E H(z)},\\
1+z&=\frac{m_\chi}{2E},\qquad 0<E\leq\frac{m_\chi}{2}.
\end{aligned}
\label{eq:cosmological-flux}
\end{equation}
Here $m_\chi$ is the parent mass, $\tau$ its lifetime, and $f$ its present dark matter fraction including the branching probability into two ALPs. The expansion rate $H(z)$ is evaluated in a flat cosmology containing matter and a cosmological constant, with a Hubble constant of 67.4~km~s$^{-1}$~Mpc$^{-1}$, matter density parameter 0.315, and present dark matter density $\rho_{\mathrm{DM},0}$ of 1.26~keV~cm$^{-3}$\cite{planck2018parameters}. The speed of light $c$ is restored in this flux expression.

% Continuum estimate: integrate Eq. (3) from 1 to 3.5 keV using
% H(z) = H0 * sqrt(0.315*(1+z)^3 + 0.685), c = 2.99792458e10 cm/s,
% 1 Mpc = 3.085677581491367e24 cm, and 1 Gyr = 3.16e16 s.
% Multiply by P = 4.4817e-18, A*Omega = 12*0.0955 cm^2 sr,
% efficiency = 0.5, and T = 23400*187 s. This is an illustrative
% constant-response continuum estimate, separate from the handout's line.
The benchmark adopts a coupling of $4.7\times10^{-11}$~GeV$^{-1}$ and a fraction divided by lifetime of $4.0\times10^{-3}$~Gyr$^{-1}$, near the selected stellar and cosmological bounds\cite{dolan2022globular,nygaard2021decaying}. A coherent field integral of 91~T~m, 12~cm$^2$ area, 0.0955~sr aperture, unit quantum efficiency, 50\% usable area, and 23,400~s per day for 187~days give approximately $4\times10^{-4}$ continuum photons between 1 and 3.5~keV. The calculation adopts these values as an illustrative exposure. The handout's separate Galactic-line estimate gives $4.1\times10^{-3}$ photons for the same decay model\cite{darkness2026handout}. Both estimates support a negative result for this benchmark. Other cosmological production models require their own intensity normalization.

\subsection{The Solar Channel}
The solar reference assumes thermal Primakoff production in the Sun through the same photon coupling that controls geomagnetic conversion. At Earth's mean distance from the Sun, the spectrum integrated over the solar disk is\cite{davoudiasl2008gecosax}
\begin{equation}
\frac{d\Phi_{a,\odot}}{dE}=6.02\times10^{10}g_{10}^{2}
 \frac{(E/\mathrm{keV})^{2.481}e^{-E/(1.205\,\mathrm{keV})}}
 {\mathrm{cm}^{2}\,\mathrm{s}\,\mathrm{keV}}.
\label{eq:solar-flux}
\end{equation}
The mean energy is 4.2~keV. In the occulted-Sun concept, ALPs traverse Earth and convert along the transparent path between the nightside atmosphere and the satellite. The instrument points toward the occulted solar core while Earth blocks direct solar X-rays. Solar production and geomagnetic conversion each depend quadratically on the coupling, giving a fourth-power dependence of the photon rate. Atmospheric transmission, field geometry, and coherence must be evaluated along this shorter path.

The published solar study projects a two-sigma coupling reach of $(4.7$--$6.6)\times10^{-11}$~GeV$^{-1}$ for masses below $10^{-4}$~eV, using approximately 1,000~cm$^2$ and $10^6$~s\cite{davoudiasl2008gecosax}. This motivates further study. DarkNESS's 12~cm$^2$ area and wide aperture require a separate forecast under the adopted stellar bound and the CAST 95\% limit of $5.8\times10^{-11}$~GeV$^{-1}$ for masses below approximately 0.02~eV\cite{cast2024extended}.

\subsection{Implication for Searches}
Diffuse counts depend on collecting area and accepted solid angle. Solar counts depend on the admitted solar flux and collecting area. A wider aperture admits additional background without increasing the solar flux once the source is enclosed. Each forecast must fold conversion through detector efficiency and usable exposure, then fit the energy and orbital dependence jointly with X-ray and particle backgrounds. Matched control exposures and particle monitoring must constrain background changes that could imitate the conversion signal.
\section{Mission Concept}
% TODO: CubeSat Mission Concept Study

\section{Results and Discussion}
% TODO: Present results and discuss their implications.

% Example figure (uncomment after adding figures/example-figure.png):
% \begin{figure}[t]
%     \centering
%     \includegraphics[width=\linewidth]{example-figure.png}
%     \caption{Describe the figure.}
%     \label{fig:example}
% \end{figure}

\section{Conclusion}
% TODO: Summarize the findings and identify next steps.

\bibliographystyle{IEEEtran}
\bibliography{references}

\end{document}
